\documentclass[12pt,a4paper]{article}
\usepackage[utf8]{inputenc}
\usepackage[spanish,es-tabla]{babel}
\usepackage{geometry}
\usepackage{graphicx}
\usepackage{multirow}
\usepackage{tabularx}
\usepackage{mathptmx}
\usepackage{fancyhdr}
\usepackage{array}
\usepackage[table]{xcolor}
\usepackage{titlesec}
\usepackage{setspace}
\usepackage{microtype}
\usepackage{lastpage}
\usepackage{booktabs}
\usepackage{enumitem}
\usepackage{amsmath}
\usepackage{amsfonts}
\usepackage{amssymb}

\usepackage{float}
\usepackage{placeins}
\usepackage{tikz}
\usepackage{pgfplots}
\pgfplotsset{compat=1.17}
\usetikzlibrary{arrows.meta,positioning,shapes.geometric,calc}

\graphicspath{{images/}{../images/}{./}}

\geometry{
    top=4.7cm,
    bottom=2.5cm,
    left=2.5cm,
    right=2.5cm,
    headheight=3.9cm,
    headsep=0.4cm
}

\setstretch{1.15}
\setlength{\parindent}{0pt}
\setlength{\parskip}{0.45em}
\setlength{\tabcolsep}{4pt}

\titleformat{\section}
  {\normalfont\fontsize{12}{14.5}\bfseries}{\thesection.}{0.6em}{\uppercase}
\titlespacing*{\section}
  {0pt}{1.8ex plus 0.4ex minus 0.2ex}{0.6ex plus 0.2ex}

\titleformat{\subsection}
  {\normalfont\fontsize{11.5}{13.5}\bfseries}{\thesubsection.}{0.5em}{}
\titlespacing*{\subsection}
  {0pt}{1.4ex plus 0.3ex minus 0.1ex}{0.4ex plus 0.1ex}

\titleformat{\subsubsection}
  {\normalfont\fontsize{10.5}{12.5}\bfseries\itshape}{\thesubsubsection.}{0.4em}{}
\titlespacing*{\subsubsection}
  {0pt}{1.1ex plus 0.2ex minus 0.1ex}{0.3ex plus 0.1ex}

\newcolumntype{Y}{>{\centering\arraybackslash}X}
\newcolumntype{L}{>{\raggedright\arraybackslash}X}
\newcolumntype{C}{>{\centering\arraybackslash}X}

\definecolor{headerblue}{RGB}{26, 54, 93}
\definecolor{darkblue}{RGB}{43, 108, 176}
\definecolor{lightgraybg}{RGB}{248, 248, 248}
\definecolor{tableheader}{RGB}{232, 238, 246}
\definecolor{navyheader}{RGB}{0, 32, 96}
\definecolor{lineblue}{RGB}{0, 102, 204}
\definecolor{areablue}{RGB}{220, 236, 252}
\definecolor{lcsred}{RGB}{217, 56, 41}
\definecolor{lcigreen}{RGB}{46, 125, 50}
\definecolor{lciorange}{RGB}{230, 81, 0}
\definecolor{ishikawared}{RGB}{160, 0, 0}

\tikzset{
    catbox/.style={draw, fill=tableheader, font=\footnotesize\bfseries, minimum width=2.8cm, minimum height=0.65cm, align=center},
    efectobox/.style={draw=ishikawared!90!black, line width=1.2pt, fill=ishikawared, text=white, font=\small\bfseries, align=center, inner sep=6pt, text width=3.8cm},
    causetext/.style={font=\scriptsize, align=right, text=black!90}
}

\newcommand{\docCodigo}{TI-SIGE-ENT-F1-005}
\newcommand{\docVersion}{1.0}
\newcommand{\docProceso}{INGENIERÍA DE SOFTWARE / COMPUERTAS LÓGICAS}
\newcommand{\docElaboracion}{24/08/2026}
\newcommand{\docRevision}{24/08/2026}
\newcommand{\docTitulo}{REGLAS DE NEGOCIO Y VALIDATION GATES LÓGICAS}
\newcommand{\docOrganizacion}{FACULTAD DE INGENIERÍA EN SISTEMAS, ELECTRÓNICA E INDUSTRIAL}
\newcommand{\docSubOrganizacion}{CARRERA DE INGENIERÍA DE SOFTWARE}

\newcommand{\encabezadofrecuente}{%
    \renewcommand{\arraystretch}{1.15}%
    \noindent\makebox[\textwidth][c]{%
    \begin{tabularx}{\textwidth}{|>{\centering\arraybackslash}m{2.8cm}|Y|>{\centering\arraybackslash}m{4.8cm}|}
        \hline
        \multirow{4}{2.8cm}{\centering\raisebox{-1.15cm}{\includegraphics[width=2.5cm,height=2.0cm,keepaspectratio]{images/logo.png}}}
        & \textbf{\fontsize{9}{11}\selectfont \docOrganizacion} 
        & \textbf{\fontsize{8}{10}\selectfont CÓDIGO:} \fontsize{8}{10}\selectfont \docCodigo \\
        \cline{3-3}
        & \textbf{\fontsize{8.5}{10.5}\selectfont \docSubOrganizacion} 
        & \textbf{\fontsize{8}{10}\selectfont VERSIÓN:} \fontsize{8}{10}\selectfont \docVersion \\
        \cline{2-3}
        & \textbf{\fontsize{8}{10}\selectfont PROCESO:} \fontsize{8}{10}\selectfont \docProceso 
        & \textbf{\fontsize{8}{10}\selectfont ELABORACIÓN:} \fontsize{8}{10}\selectfont \docElaboracion \\
        \cline{2-3}
        & \textbf{\fontsize{8.5}{10.5}\selectfont \docTitulo} 
        & \textbf{\fontsize{8}{10}\selectfont ÚLTIMA REVISIÓN:} \fontsize{8}{10}\selectfont \docRevision \\
        \hline
    \end{tabularx}%
    }%
}

\pagestyle{fancy}
\fancyhf{}
\fancyhead[C]{\encabezadofrecuente}
\fancyfoot[C]{\fontsize{10}{12}\selectfont Página \thepage\ de \pageref{LastPage}}
\renewcommand{\headrulewidth}{0pt}

\begin{document}

\noindent{\fontsize{10.5}{12.5}\selectfont \textbf{Integrantes:} Heidi Villavicencio \quad$\vert$\quad Steeven Toala \quad$\vert$\quad Nixon Hurtado \quad$\vert$\quad Andrea Vásquez}

\vspace{0.25cm}

\section{Título}
\textbf{Denominación Oficial del Entregable:} Reglas de Negocio y Validation Gates Lógicas -- Código EDT: 1.5

\section{Introducción}
El problema principal que degrada los procesos documentales del SGC radica en que los controles de calidad se efectúan tardíamente, tras semanas de reproceso. La arquitectura de software moderna soluciona este cuello de botella mediante compuertas de validación (\textit{Validation Gates}) que operan como middlewares lógicos en el backend, impidiendo que transiciones erróneas o estados inconsistentes se persistan en la base de datos. Este entregable formaliza las 10 Reglas de Negocio (RN-01 a RN-10) y sus correspondientes compuertas de validación que blindan la integridad del sistema SIGE.

\section{Objetivo}
\subsection{Objetivo General}
Especificar formal y matemáticamente las 10 reglas de negocio institucionales y diseñar las compuertas lógicas preventivas (\textit{Validation Gates}) que controlan las transiciones documentales en el backend de SIGE.
\subsection{Objetivos Específicos}
\begin{itemize}[leftmargin=1.5em, itemsep=2pt]
    \item Modelar las condiciones matemáticas de sumatoria y consistencia horaria en las matrices de actividades.
    \item Definir las precondiciones y códigos HTTP de bloqueo semántico (400, 422, 403) para cada compuerta.
    \item Establecer la regla de cascada de firmas que invalida automáticamente firmas previas ante correcciones.
\end{itemize}

\section{Alcance}
\subsection{Límites del Entregable (In-Scope)}
\begin{itemize}[leftmargin=1.5em, itemsep=2pt]
    \item 10 Reglas de Negocio institucionales (RN-01 a RN-10).
    \item Compuertas: Gate 1 (Borrador), Gate 2 (Envío a Revisión), Gate 3 (Aval Coordinación) y Gate 4 (Aprobación Decanato).
\end{itemize}
\subsection{Exclusiones (Out-of-Scope)}
\begin{itemize}[leftmargin=1.5em, itemsep=2pt]
    \item Reglas contables financieras de nómina docente.
\end{itemize}

\section{Definiciones}
\begin{itemize}[leftmargin=1.5em, itemsep=2pt]
    \item \textbf{Validation Gate:} Mecanismo interceptor que evalúa invariantes del dominio antes de autorizar la persistencia o cambio de estado.
    \item \textbf{Invariante de Ponderación:} Condición estricta que exige que la suma porcentual de actividades sea exactamente igual a $100\%$.
    \item \textbf{Inmutabilidad de Aprobación:} Estado de congelamiento en el cual un documento aprobado no puede ser modificado por ningún rol sin autorización explícita de apertura extemporánea.
\end{itemize}

\section{Responsabilidades}
\begin{itemize}[leftmargin=1.5em, itemsep=2pt]
    \item \textbf{Steeven Toala / Nixon Hurtado (Responsables EDT 1.5):} Modelado matemático de compuertas lógicas y especificación técnica de middlewares.
    \item \textbf{Heidi Villavicencio (Analista):} Validación de congruencia de las reglas con la normativa SGC.
    \item \textbf{Ing. Paulo Torres (Docente Evaluador):} Aprobación de la rigurosidad de las compuertas preventivas.
\end{itemize}

\section{Normativa Legal y Técnica}
\subsection{Marco Legal Ecuatoriano}
\begin{itemize}[leftmargin=1.5em, itemsep=2pt]
    \item \textbf{Régimen Académico (CES):} Cumplimiento estricto del porcentaje de avance en horas de cátedra y actividades extracurriculares.
\end{itemize}

\subsection{Normas Técnicas Internacionales}
\begin{itemize}[leftmargin=1.5em, itemsep=2pt]
    \item \textbf{ISO 9001:2015, Cláusula 8.5.1 (Control de la Producción y Provisión del Servicio):} Implementación de acciones preventivas para evitar no conformidades.
    \item \textbf{RFC 7231 (HTTP/1.1 Semantics):} Uso riguroso de códigos de error semánticos (400 Bad Request, 422 Unprocessable Entity).
\end{itemize}

\section{Método, Políticas y Contenidos}
\subsection{Matriz Oficial de Reglas de Negocio y Validation Gates (RN-01 a RN-10)}
A continuación se detalla la matriz que rige las validaciones automáticas del sistema:

\begin{table}[H]
\centering
\footnotesize
\renewcommand{\arraystretch}{1.15}
\begin{tabularx}{\textwidth}{|c|>{\raggedright\arraybackslash\bfseries}m{2.2cm}|X|c|c|}
\hline
\rowcolor{tableheader}
\textbf{ID} & \textbf{Regla} & \textbf{Condición Lógica de Bloqueo / Validación Preventiva} & \textbf{Entidad} & \textbf{HTTP} \\
\hline
\textbf{RN-01} & Suma 100\% Plan & Bloquear envío a revisión si $\sum_{i=1}^n \text{Ponderación}_i \neq 100\%$. & PlanTrabajo & 422 \\
\hline
\textbf{RN-02} & Cronograma Válido & Toda actividad debe tener fechas inicio/fin dentro del periodo lectivo oficial. & Actividad & 400 \\
\hline
\textbf{RN-03} & Evidencia Obligatoria & Bloquear informe si avance declarado $> 0\%$ y no existen archivos adjuntos en S3. & Informe & 422 \\
\hline
\textbf{RN-04} & Justificación de Brecha & Exigir texto explicativo obligatorio si el avance es menor al planificado. & Informe & 400 \\
\hline
\textbf{RN-05} & Cascada Secuencial & No se permite firma de Coordinador sin firma previa del Docente Elaborador. & Firmas & 403 \\
\hline
\textbf{RN-06} & Firma de Autoridad & No se permite firma de Decano sin aval y firma previa del Coordinador. & Firmas & 403 \\
\hline
\textbf{RN-07} & Invalidación por Cambio & Toda modificación a un documento avalado invalida firmas previas (retorno a borrador). & Flujo & 200 \\
\hline
\textbf{RN-08} & Unicidad por Período & Un docente solo puede tener un Plan de Trabajo activo por comisión en cada ciclo. & PlanTrabajo & 409 \\
\hline
\textbf{RN-09} & Inmutabilidad de Aprobado & Una vez firmado por Decano, el registro pasa a estado de solo lectura inalterable. & Seguridad & 403 \\
\hline
\textbf{RN-10} & Extemporaneidad Auditada & Desbloqueos fuera de fecha requieren autorización y quedan en log inmutable. & Auditoría & 200 \\
\hline
\end{tabularx}
\caption{Matriz de Reglas de Negocio y Compuertas Lógicas Preventivas en SIGE.}
\label{tab:reglas_negocio}
\end{table}
\FloatBarrier

\section{Descripción de actividades}
\subsection{Formulación Matemática de las Compuertas de Validación}
La compuerta más crítica (\textbf{Gate 2}) evalúa la integridad del plan antes de transicionar de \texttt{borrador} a \texttt{en\_revision}:
\begin{equation}
\text{Gate}_2(\text{Plan}) = \left( \sum_{i=1}^{n} w_i = 100\% \right) \;\land\; \left( \forall i, \; t_{\text{inicio}, i} < t_{\text{fin}, i} \right) \;\land\; \left( n \ge 1 \right)
\end{equation}
Donde $w_i$ representa el peso ponderado de la actividad $i$, y $n$ el número de actividades registradas. Si $\text{Gate}_2$ resulta falsa, la petición REST es abortada con código 422 y un cuerpo JSON que especifica el error de negocio.

\section{Conclusiones}
\begin{itemize}[leftmargin=1.5em, itemsep=2pt]
    \item La implementación de compuertas lógicas preventivas elimina de raíz los reprocesos aritméticos que afectaban al 42\% de los documentos As-Is.
    \item La regla de invalidación en cascada garantiza la pureza legal de los documentos, impidiendo que existan firmas aprobatorias sobre textos alterados posteriormente.
\end{itemize}

\section{Firmas de Responsabilidad}

\noindent
\renewcommand{\arraystretch}{1.2}
\begin{tabularx}{\textwidth}{|>{\raggedright\arraybackslash}m{2.9cm}|>{\raggedright\arraybackslash}X|>{\centering\arraybackslash}m{4.6cm}|>{\centering\arraybackslash}m{2.3cm}|}
    \hline
    \rowcolor{tableheader}
    \textbf{ACCIÓN} & \textbf{NOMBRE Y CARGO} & \textbf{FIRMA} & \textbf{FECHA} \\
    \hline
    \textbf{Elaborado por:} & Heidi Villavicencio\newline Analista de Calidad / UI-UX & \parbox[c]{4.4cm}{\centering\vspace{4pt}\includegraphics[height=1.5cm,width=4.0cm,keepaspectratio]{images/firma_heidi.png}\vspace{4pt}} & 24/08/2026 \\
    \hline
    \textbf{Revisado por:} & Steeven Toala\newline Líder de Proyecto / Arquitecto & \parbox[c]{4.4cm}{\centering\vspace{4pt}\includegraphics[height=1.5cm,width=4.0cm,keepaspectratio]{images/firma_steveen.png}\vspace{4pt}} & 24/08/2026 \\
    \hline
    \textbf{Aprobado por:} & Ing. Paulo Torres, Mg.\newline Docente Evaluador -- Gestión Proyectos & \parbox[c]{4.4cm}{\centering\vspace{8pt}\textbf{\textcolor{gray!80}{\scriptsize [ PENDIENTE DE APROBACIÓN ]}}\vspace{8pt}} & \textbf{\textcolor{gray!90}{Pendiente}} \\
    \hline
\end{tabularx}

\vspace{0.4cm}

\section{Control de Cambios}

\noindent
\renewcommand{\arraystretch}{1.2}
\begin{tabularx}{\textwidth}{|>{\hsize=0.45\hsize\centering\arraybackslash}X|>{\hsize=1.55\hsize\raggedright\arraybackslash}X|>{\hsize=1.2\hsize\raggedright\arraybackslash}X|>{\hsize=0.8\hsize\centering\arraybackslash}X|}
    \hline
    \rowcolor{tableheader}
    \textbf{Versión} & \textbf{Descripción del cambio} & \textbf{Responsable del cambio} & \textbf{Fecha de aprobación} \\
    \hline
    1.0 & Emisión inicial del entregable según EDT del Tablero Trello & Equipo de Proyecto SIGE & 24/08/2026 \\
    \hline
\end{tabularx}

\vspace{0.4cm}

\section{Anexos}

\subsection*{Anexo 1: Ficha de Registro y Control en el Tablero Oficial de Gestión (EDT Trello)}
\begin{table}[H]
\centering
\small
\renewcommand{\arraystretch}{1.2}
\begin{tabularx}{\textwidth}{|>{\bfseries}m{3.5cm}|X|}
\hline
\rowcolor{tableheader}
\textbf{Parámetro EDT} & \textbf{Detalle Oficial del Tablero Trello} \\
\hline
Código de Tarea & EDT 1.5 \\
\hline
Nombre de la Tarjeta & Reglas de negocio y Validation Gates lógicas \\
\hline
Fase Correspondiente & FASE 1. ANÁLISIS Y PLANIFICACIÓN \\
\hline
Período de Ejecución & 19 de agosto al 24 de agosto de 2026 (4 días laborables) \\
\hline
Horas Asignadas & 32 horas hombre netas \\
\hline
Responsables Técnicos & Steeven Toala (ST), Nixon Hurtado (NH) \\
\hline
Entregable Oficial & \texttt{05\_Reglas\_Negocio\_Validation\_Gates.pdf} \\
\hline
Estado en el Tablero & 100\% Completado y Verificado en Reunión de Hito 1 \\
\hline
\end{tabularx}
\caption{Ficha técnica de la tarea EDT 1.5 en el Tablero Trello.}
\end{table}

\subsection*{Anexo 2: Evidencia de Pruebas Lógicas de la Compuerta de Validación}
Resultados de ejecución del script de validación preliminar en backend que confirma el rechazo del código 422 ante sumatorias de actividades incompletas ($85\% \to \text{RECHAZADO}$).

\subsection*{Anexo 3: Trazabilidad en Repositorio Institucional Git}
\begin{itemize}[leftmargin=1.5em, itemsep=2pt]
    \item \textbf{Repositorio:} \texttt{FISEI\_Project / SIGE}
    \item \textbf{Rama de Trabajo:} \texttt{main}
    \item \textbf{Commit de Consolidación:} \texttt{b72183c - feat(core): definicion formal de compuertas validation gates RN-01 a RN-10}
    \item \textbf{Archivos Afectados:} \texttt{src/middlewares/validationGates.js}, \texttt{Entregables\_Latex/Fase\_1/1.5*.tex}
\end{itemize}


\end{document}
